% Part: applied-modal-logics
% Chapter: epistemic-logic
% Section: public-announcement-logic-lang

\documentclass[../../../include/open-logic-section]{subfiles}

\begin{document}

\olfileid{aml}{el}{pal}

\olsection{منطق الإعلان العام}

تمثل المنطقيات المعرفية الديناميكية تغير معرفة الفاعلين بمرور الزمن أو
بورود معلومات جديدة. وأكثرها يصور هذا التغير بــ\emph{أحداث} أو
\emph{تحديثات} معلوماتية. وأبسط التحديث إعلان عام: تعلن فيه صيغة إعلانًا
صادقًا ويشهد الإعلان جميع الفاعلين معًا. ولأجل ذلك نوسع اللغة كما يأتي:

\begin{defn}
لتكن ${\OLId{latin-looped}{G}}$ مجموعة رموز الفاعلين. فاللغة الأساسية لمنطق المعرفة المتعدد
الفاعلين ذي الإعلانات العامة تتركب مما يأتي:
\begin{enumerate}
  \tagitem{prvFalse}{الثابت القضي لل!!{falsity}~$\lfalse$.}{}
  \tagitem{prvTrue}{الثابت القضي لل!!{truth}~$\ltrue$.}{}
  \item مجموعة !!{denumerable} من ال!!{propositional variable}s:
    $\Obj p_{\OLMathNumeral{0}}$، $\Obj p_{\OLMathNumeral{1}}$، $\Obj p_{\OLMathNumeral{2}}$، \dots
  \item الروابط القضوية: \startycommalist
  \iftag{prvNot}{\ycomma $\lnot$ (النفي)}{}%
  \iftag{prvAnd}{\ycomma $\land$ (العطف)}{}%
  \iftag{prvOr}{\ycomma $\lor$ (الفصل)}{}%
  \iftag{prvIf}{\ycomma $\lif$ (!!{conditional})}{}%
  \iftag{prvIff}{\ycomma $\liff$ (!!{biconditional})}{}%
  \item مؤثر المعرفة $\Knows_{\OLId{latin-ordinary}{a}}$، حيث ${\OLId{latin-ordinary}{a}} \in {\OLId{latin-looped}{G}}$.
  \item مؤثر الإعلان العام $[!B]$، حيث $!B$ هي !!a{formula}.
\end{enumerate}
\end{defn}

ولما كان مؤثر الإعلان العام جاريا مجرى مؤثر الصندوق، كان حد اللغة
الاستقرائي على الوجه الآتي:

\begin{defn}
  تحد \emph{!!^{formula}s} اللغة المعرفية استقرائيًا كما يأتي:
  \begin{enumerate}
  \tagitem{prvFalse}{$\lfalse$ !!{formula} ذرية.}{}

  \tagitem{prvTrue}{$\ltrue$ !!{formula} ذرية.}{}

  \item كل متغير قضي $\Obj p_{\OLId{latin-ordinary}{i}}$ هو !!{formula} (ذرية).

  \tagitem{prvNot}{إذا كانت $!A$ !!a{formula}، فإن $\lnot !A$
  !!a{formula}.}{}

  \tagitem{prvAnd}{إذا كانت $!A$ و$!B$ !!{formula}s، فإن $(!A \land
  !B)$ !!a{formula}.}{}

  \tagitem{prvOr}{إذا كانت $!A$ و$!B$ !!{formula}s، فإن $(!A \lor !B)$
  !!a{formula}.}{}

  \tagitem{prvIf}{إذا كانت $!A$ و$!B$ !!{formula}s، فإن $(!A \lif !B)$
  !!a{formula}.}{}

  \tagitem{prvIff}{إذا كانت $!A$ و$!B$ !!{formula}s، فإن $(!A \liff !B)$
  !!a{formula}.}{}

  \item إذا كانت $!A$ !!a{formula} وكان ${\OLId{latin-ordinary}{a}} \in {\OLId{latin-looped}{G}}$، فإن $\Knows_{\OLId{latin-ordinary}{a}} !A$
  !!a{formula}.
  
  \item إذا كانت $!A$ و$!B$ !!{formula}s، فإن $[!A] !B$ !!a{formula}.

  \tagitem{limitClause}{ولا شيء غير ذلك !!a{formula}.}{}
\end{enumerate}
\end{defn}

والمقصود بال!!{formula} $[!A] !B$ قولنا: «بعد أن تُعلن $!A$ إعلانًا
صادقًا، تصدق~$!B$.» وقد نحتاج في باب الإعلانات العامة إلى ذكر المعرفة
المشتركة أيضًا؛ فحينئذ تشتمل اللغة على مؤثرها~$\CKnows_{\OLId{latin-looped}{G}} !A$.

\end{document}
